%! Function Model Selection and Assumption Articulation — Unit 1, Lesson 1.13 — Homework (Answer Key)
\documentclass[10pt]{article}
\usepackage{precalculus-article}
\usepackage{precalculus-key}
\newcommand{\TallMath}[1]{$\displaystyle #1\rule[-1.4em]{0pt}{3.2em}$}

\begin{document}

\pageheader{Unit 1, Lesson 1.13}{Homework --- Answer Key}
\namedateperiod

\vspace{0.1in}

\begin{enumerate}[itemsep=14pt]

\item \textbf{Difference table.}

      \medskip
      \renewcommand{\arraystretch}{1.4}
      \begin{tabular}{|c|c|c|c|c|}
      \hline
      $x$ & $y$ & $\Delta y$ & $\Delta^2 y$ & $\Delta^3 y$ \\
      \hline
      0 & 2   & ---         & ---         & ---         \\
      \hline
      1 & 6   & \ans{4} & ---      & ---         \\
      \hline
      2 & 16  & \ans{10} & \ans{6} & ---    \\
      \hline
      3 & 38  & \ans{22} & \ans{12} & \ans{6} \\
      \hline
      4 & 78  & \ans{40} & \ans{18} & \ans{6} \\
      \hline
      5 & 142 & \ans{64} & \ans{24} & \ans{6} \\
      \hline
      \end{tabular}

      \medskip
      The \ans{third} differences are constant, so this data is best modeled by a
      degree-\ans{3} polynomial.

\item \textbf{Context-to-model matching.}

      \medskip
      A $\to$ \ans{3} \qquad B $\to$ \ans{1} \qquad C $\to$ \ans{2} \qquad D $\to$ \ans{1}

      \begin{teachernote}
      B and D are both linear, but they match to model type 1 independently.
      Profit per unit being constant (D) is also linear. C is quadratic because projectile
      height as a function of horizontal distance follows $h = ax^2 + bx + c$.
      \end{teachernote}

\item \textbf{Assumption articulation.}

      (a) \ansline{This model assumes the cyclist travels at a constant (unchanging) speed throughout the trip.}

      (b) Domain: \ans{$t \geq 0$} \quad Reason: \ansline{Time cannot be negative; the cyclist starts at $t = 0$.}

\item \textbf{Piecewise model description.}

      Piece 1 (when \ans{$0 < h \leq 1$}): $C = $ \ans{$3.00$} \quad Type: \ans{constant}

      Piece 2 (when \ans{$1 < h \leq k$, where $3 + 1.5(h-1) < 15$}): $C = $ \ans{$3.00 + 1.50(h - 1)$} \quad Type: \ans{linear}

      Piece 3 (when \ans{$h > k$ for the maximum threshold}): $C = $ \ans{$15.00$} \quad Type: \ans{constant}

      \begin{teachernote}
      The daily cap is reached when $3 + 1.5(h-1) = 15 \Rightarrow h = 9$ hours.
      So Piece 2 applies for $1 < h \leq 9$ and Piece 3 for $h > 9$. Accept reasonable interval notation.
      \end{teachernote}

\item \textbf{Model selection with justification.}

      \medskip
      \renewcommand{\arraystretch}{1.3}
      \begin{tabular}{|c|c|c|c|}
      \hline
      $r$ & $A$ & $\Delta A$ & $\Delta^2 A$ \\
      \hline
      1 & 3   & ---         & ---         \\
      \hline
      2 & 12  & \ans{9} & ---      \\
      \hline
      3 & 27  & \ans{15} & \ans{6} \\
      \hline
      4 & 48  & \ans{21} & \ans{6} \\
      \hline
      \end{tabular}

      \medskip
      Model type: \ans{Quadratic}

      Context clue: \ansline{Area is a 2-D (flat surface) measurement, which matches the area $\propto r^2$ shape of a circle or spreading region.}

\end{enumerate}

\vspace{0.1in}

\begin{extensionbox}
\textbf{Extension (optional).}

\medskip
\renewcommand{\arraystretch}{1.3}
\begin{tabular}{|c|c|c|c|c|}
\hline
$x$ & $y$ & $\Delta y$ & $\Delta^2 y$ & $\Delta^3 y$ \\
\hline
0 & 0  & ---         & ---         & ---         \\
\hline
1 & 2  & \ans{2} & ---      & ---         \\
\hline
2 & 8  & \ans{6} & \ans{4} & ---    \\
\hline
3 & 18 & \ans{10} & \ans{4} & \ans{0} \\
\hline
\end{tabular}

\medskip
\ansline{The student's claim is partially correct but imprecise. With 4 points, a degree-3 polynomial is the maximum needed; however, the second differences are already constant (all 4), so the data is quadratic (degree 2). A cubic is not required. The unique polynomial of \textit{lowest} degree through these 4 points is actually quadratic ($y = 2x^2$).}
\end{extensionbox}

\end{document}
